\documentclass[12pt]{article}

% --- Idioma ---
\usepackage[spanish,es-nodecimaldot]{babel}

% --- Fuente Times New Roman (compilar con XeLaTeX) ---
\usepackage{fontspec}
\setmainfont{Times New Roman}

% --- Márgenes 2.5 cm por lado ---
\usepackage[top=2.5cm, bottom=2.5cm, left=2.5cm, right=2.5cm]{geometry}

% --- Interlineado 1.5 ---
\usepackage{setspace}
\onehalfspacing

% --- Espaciado posterior 6 pt, sin sangría de párrafo ---
\setlength{\parindent}{0pt}
\setlength{\parskip}{6pt}

% --- Sangría francesa en referencias (APA) ---
\usepackage{hanging}

% --- Hipervínculos ---
\usepackage[hidelinks]{hyperref}

% ============================================================
\begin{document}
% ============================================================

\section{Marco Teórico}

\subsection{La salud como capacidad fundamental y dimensión de la pobreza multidimensional}

El análisis económico del sector salud requiere partir de una distinción conceptual previa: la salud
no es únicamente un bien de consumo ni un insumo productivo, sino una capacidad fundamental del
ser humano. Esta distinción tiene implicancias directas sobre cómo se mide el bienestar y sobre qué
papel le corresponde al Estado en la provisión de servicios sanitarios.

Sen (1988) argumenta que la enfermedad no es simplemente un síntoma observable de pobreza,
sino una restricción a las capacidades básicas de las personas: priva al individuo de su libertad para
funcionar, trabajar y participar en la vida social. Este argumento es ampliado en Sen (2000), donde
el desarrollo económico es reinterpretado no como crecimiento del producto per cápita, sino como
expansión de las libertades reales que los individuos pueden ejercer. Bajo este enfoque, un país que
duplica su PIB sin reducir la mortalidad infantil o aumentar la esperanza de vida no puede
considerarse plenamente desarrollado.

La operacionalización de este enfoque es realizada por Alkire y Foster (2011), quienes construyen
el Índice de Pobreza Multidimensional (IPM) a partir de tres dimensiones: salud, educación y nivel
de vida. La dimensión salud se mide a través de indicadores como la mortalidad infantil y la
desnutrición crónica. Su método identifica como pobre multidimensional a quien carece del umbral
mínimo en al menos un tercio de los indicadores ponderados, lo que permite capturar la intensidad
de la pobreza más allá del ingreso monetario. Esta metodología es también la base del Índice de
Desarrollo Humano (IDH) del Programa de las Naciones Unidas para el Desarrollo, que incorpora
la esperanza de vida al nacer como \textit{proxy} de la dimensión salud. Para sintetizar indicadores
heterogéneos en una medida única, el IDH recurre a técnicas de índices compuestos, entre ellas el
Análisis de Componentes Principales (PCA), que permite reducir la dimensionalidad de un conjunto
de variables correlacionadas preservando la mayor varianza posible (UNDP, 2015).

La relevancia empírica de esta distinción se evidencia en el caso peruano. El INEI (2016) reporta
que la pobreza monetaria cayó de 58.7\,\% en 2004 a 21.8\,\% en 2015, impulsada por el
crecimiento económico sostenido. Sin embargo, como señala Tezano (2013), la reducción de la
pobreza multidimensional fue considerablemente más lenta, en particular en las regiones de sierra y
selva, donde las brechas en acceso a servicios de salud y saneamiento persistieron incluso durante
los años de mayor expansión del PIB. Esta divergencia sugiere que el crecimiento es condición
necesaria pero no suficiente para mejorar el bienestar en salud, lo que motiva el análisis de la
relación entre gasto público y pobreza que se desarrolla en la sección de series de tiempo.

\subsection{La salud como capital humano: efectos sobre la productividad y el crecimiento}

La segunda línea teórica relevante aborda la salud desde la perspectiva del capital humano. La
teoría del capital humano postula que los individuos invierten en educación y salud con el fin de
incrementar su productividad futura (Becker, 1964). Bajo este marco, el gasto en salud no es un
gasto corriente de consumo, sino una inversión que aumenta el stock de capital humano de la
economía y, por tanto, su capacidad productiva de largo plazo. A nivel macroeconómico, esta idea
se formaliza en las extensiones del modelo de Solow que incorporan capital humano como factor de
producción: una mejora sostenida en el estado de salud de la fuerza laboral eleva la productividad
total de los factores y desplaza hacia arriba la trayectoria de crecimiento de estado estacionario
(De~Gregorio, 2007). Esto implica que el gasto público en salud tiene un efecto potencialmente
positivo sobre el crecimiento del PIB, no solo a través del canal de demanda agregada de corto
plazo, sino a través de la acumulación de capital humano en el largo plazo.

La evidencia para el Perú es analizada empíricamente por Cortez (1999), quien estima el efecto de
las condiciones de salud sobre la productividad laboral e ingresos individuales. Sus resultados
muestran retornos significativos, con diferencias importantes por género y región geográfica: los
efectos son más pronunciados en zonas rurales de sierra y selva, precisamente donde la cobertura de
salud pública es más precaria y la carga de enfermedad es mayor. Valdez et al. (2013) complementan
este diagnóstico con el perfil epidemiológico del Perú: las enfermedades no transmisibles
---cardiovasculares, diabetes, cáncer--- y las enfermedades diarreicas agudas concentran la mayor
carga de morbilidad y constituyen los principales factores que reducen la productividad laboral. Este
hallazgo es consistente con la relación empírica entre acceso a agua potable segura y mortalidad
infantil que se estima en la sección econométrica, dado que el saneamiento básico es el principal
determinante de las enfermedades diarreicas agudas en el país.

\subsection{Fallas de mercado y el rol del Estado en la provisión de salud}

La justificación teórica de la intervención pública en el sector salud descansa en la presencia de
fallas de mercado que impiden que el equilibrio privado sea eficiente. La primera y más relevante es
la asimetría de información: el paciente no puede evaluar \textit{ex ante} la calidad del servicio
médico ni el precio justo del medicamento, lo que genera una relación de agencia entre el médico y
el paciente y abre espacio para comportamientos oportunistas en la cadena de comercialización
farmacéutica. Esta falla justifica la regulación de precios de medicamentos y la fiscalización por
parte de DIGEMID, cuya función es reducir las rentas informacionales que capturan las cadenas de
farmacias en perjuicio del consumidor final.

La segunda falla son las externalidades positivas: la vacunación y el acceso al saneamiento básico
generan beneficios sociales que exceden ampliamente los beneficios privados, por lo que el mercado
los subprovee en ausencia de intervención pública. La tercera falla es el carácter de bien meritorio
de la salud: la sociedad considera que el acceso a servicios sanitarios básicos debe garantizarse con
independencia de la capacidad de pago del individuo, lo que justifica la existencia del Seguro
Integral de Salud (SIS) y de EsSalud como mecanismos de aseguramiento colectivo. Estas tres
fallas ---asimetría de información, externalidades positivas y bien meritorio--- configuran el marco
dentro del cual opera el sector salud peruano y explican las tensiones entre el sistema público y el
privado que se analizan en las secciones de composición e institucionalidad de este trabajo.

% ============================================================
\section*{Referencias}
% ============================================================

\begin{hangparas}{1.5em}{1}

Alkire, S., \& Foster, J. (2011). Counting and multidimensional poverty measurement.
\textit{Journal of Public Economics, 95}(7), 476--487.

Becker, G. (1964). \textit{Human Capital: A Theoretical and Empirical Analysis, with Special
Reference to Education}. New York: National Bureau of Economic Research.

Cortez, R. (1999). \textit{Salud y productividad en el Perú: Un análisis empírico por género y
región} (Working Paper R-363). Washington D.C.: Inter-American Development Bank.

De Gregorio, J. (2007). \textit{Macroeconomía: Teoría y Políticas}. Santiago: Pearson
Educación.

Instituto Nacional de Estadística e Informática. (2016). \textit{Evolución de la Pobreza
Monetaria 2009-2015. Informe Técnico}. Lima: INEI.

Sen, A. (1988). Property and hunger. \textit{Economics and Philosophy, 4}(1), 57--68.

Sen, A. (2000). \textit{Desarrollo y Libertad}. Buenos Aires: Planeta.

Tezano, S. (2013). \textit{Desarrollo humano, pobreza y desigualdades}. Santander: Universidad
de Cantabria, Cátedra de Cooperación Internacional y con Iberoamérica.

United Nations Development Programme. (2015). \textit{Human Development Report 2015:
Technical Notes --- Calculating the Indices}. New York: UNDP. Recuperado de
\url{http://hdr.undp.org/en/}

Valdez, W., Napanga, E., Oyola, A., Mariños, J., Vílchez, A., Medina, J., \& Berto, M. (2013).
\textit{Análisis de la Situación de Salud en el Perú}. Lima: Ministerio de Salud.

\end{hangparas}

\end{document}
